
% =========================================================== 8. Caso Brasil, os dados
\begin{frame}
  \framesubtitle{Brazil case \ \textbullet\ \ Wednesday, 4 February 2026}
  \frametitle{The fleet climbs 43 GW while MC multiplies by 8.}

  \begin{columns}[T,onlytextwidth]
  \begin{column}{0.545\textwidth}
    \vspace{0.5em}
    \centering
    \resizebox{\linewidth}{!}{\input{figures/fig-dia-en.tikz}}

    \vspace{0.5em}
    \raggedright
    {\fontsize{5.9}{7.4}\selectfont\color{muted}
     \textcolor{ink}{\rule[0.22ex]{0.9em}{1.1pt}}~net load\hfill
     \textcolor{muted}{\rule[0.22ex]{0.9em}{0.5pt}}~gross load\hfill
     \textcolor{tealx}{\rule[0.22ex]{0.9em}{0.7pt}}~wind + solar\hfill
     \textcolor{brick!35}{\rule[0ex]{0.6em}{0.6em}}~MC, right axis\par}

    \vspace{0.4em}
    \gloss{t1 is the trough (12:00) and t2 the peak (19:00) of the controllable
    fleet curve. Gross load barely moves: what climbs 43 GW is the fleet.}
  \end{column}

  \begin{column}{0.425\textwidth}
    \vspace{0.5em}
    {\fontsize{7.3}{9.2}\selectfont
    \setlength{\tabcolsep}{2.8pt}
    \renewcommand{\arraystretch}{1.12}
    \begin{tabular}{l r r r}
      \toprule
       & \textbf{t1 · 12:00} & \textbf{t2 · 19:00} & \textbf{$\Delta$} \\
      \midrule
      Gross load         & 97,168 & 99,352 & \hl{$+$2,185} \\
      Wind + solar       & 41,281 & 5,464  & \hl{$-$35,816} \\
      Distributed        & 28,428 & 278    & \hl{$-$28,150} \\
      \addlinespace[0.1em]\midrule
      Hydro              & 41,696 & 81,048 & $+$39,352 \\
      Thermal            & 8,453  & 12,064 & $+$3,611 \\
      \textbf{Controllable} & \hg{50,148} & \hg{93,111} & \hg{$+$42,963} \\
      \addlinespace[0.1em]\midrule
      MC (SE)            & 309    & \hl{2,521} & $\times$8.2 \\
      Spot (SE)          & 350    & 1,200  & $\times$3.4 \\
      \bottomrule
    \end{tabular}}

    \vspace{0.5em}
    \thesis{\fontsize{7.1}{9}\selectfont
      Two days later the same system moved 27 GW and MC moved 47 R\$/MWh. On the
      4th it moved 43 GW and MC moved 2,212. The value of flexibility is
      concentrated in a handful of hours of the year.}
  \end{column}
  \end{columns}
\end{frame}

% =========================================================== 9. Caso Brasil, a pilha
\begin{frame}
  \framesubtitle{Brazil case \ \textbullet\ \ the stack and the battery}
  \frametitle{The hydro ramp is what separates the two prices.}

  \begin{columns}[T,onlytextwidth]
  \begin{column}{0.545\textwidth}
    \vspace{0.5em}
    \centering
    \resizebox{\linewidth}{!}{\begin{tikzpicture}[x=1cm,y=1cm]
  \fill[muted,opacity=.42] (0.000,0) rectangle (2.860,0.100);
  \fill[brick,opacity=.26] (2.860,0) rectangle (2.863,0.000);
  \fill[brick,opacity=.26] (2.863,0) rectangle (2.867,0.000);
  \fill[brick,opacity=.26] (2.867,0) rectangle (2.876,0.000);
  \fill[brick,opacity=.26] (2.876,0) rectangle (2.888,0.000);
  \fill[brick,opacity=.26] (2.888,0) rectangle (2.900,0.000);
  \fill[brick,opacity=.26] (2.900,0) rectangle (2.924,0.000);
  \fill[brick,opacity=.26] (2.924,0) rectangle (2.968,0.000);
  \fill[brick,opacity=.26] (2.968,0) rectangle (2.999,0.058);
  \fill[brick,opacity=.26] (2.999,0) rectangle (3.023,0.197);
  \fill[brick,opacity=.26] (3.023,0) rectangle (3.059,0.216);
  \fill[brick,opacity=.26] (3.059,0) rectangle (3.098,0.238);
  \fill[brick,opacity=.26] (3.098,0) rectangle (3.103,0.277);
  \fill[brick,opacity=.26] (3.103,0) rectangle (3.108,0.277);
  \fill[brick,opacity=.26] (3.108,0) rectangle (3.113,0.277);
  \fill[brick,opacity=.26] (3.113,0) rectangle (3.119,0.277);
  \fill[brick,opacity=.26] (3.119,0) rectangle (3.125,0.277);
  \fill[brick,opacity=.26] (3.125,0) rectangle (3.137,0.277);
  \fill[brick,opacity=.26] (3.137,0) rectangle (3.177,0.277);
  \fill[brick,opacity=.26] (3.177,0) rectangle (3.179,0.424);
  \fill[brick,opacity=.26] (3.179,0) rectangle (3.180,0.504);
  \fill[brick,opacity=.26] (3.180,0) rectangle (3.189,0.537);
  \fill[brick,opacity=.26] (3.189,0) rectangle (3.212,0.546);
  \fill[brick,opacity=.26] (3.212,0) rectangle (3.235,0.546);
  \fill[brick,opacity=.26] (3.235,0) rectangle (3.344,0.558);
  \fill[brick,opacity=.26] (3.344,0) rectangle (3.394,0.569);
  \fill[brick,opacity=.26] (3.394,0) rectangle (3.412,0.573);
  \fill[tealx,opacity=.28] (3.412,0) rectangle (6.111,0.573);
  \fill[brick,opacity=.26] (6.111,0) rectangle (6.136,0.575);
  \fill[brick,opacity=.26] (6.136,0) rectangle (6.153,0.576);
  \fill[brick,opacity=.26] (6.153,0) rectangle (6.184,0.585);
  \fill[brick,opacity=.26] (6.184,0) rectangle (6.209,0.590);
  \fill[brick,opacity=.26] (6.209,0) rectangle (6.226,0.593);
  \fill[brick,opacity=.26] (6.226,0) rectangle (6.244,0.636);
  \fill[brick,opacity=.26] (6.244,0) rectangle (6.266,0.645);
  \fill[brick,opacity=.26] (6.266,0) rectangle (6.281,0.755);
  \fill[brick,opacity=.26] (6.281,0) rectangle (6.288,0.773);
  \fill[brick,opacity=.26] (6.288,0) rectangle (6.402,0.828);
  \fill[brick,opacity=.26] (6.402,0) rectangle (6.410,0.837);
  \fill[brick,opacity=.26] (6.410,0) rectangle (6.417,0.964);
  \fill[brick,opacity=.26] (6.417,0) rectangle (6.441,0.965);
  \fill[brick,opacity=.26] (6.441,0) rectangle (6.532,1.146);
  \fill[brick,opacity=.26] (6.532,0) rectangle (6.540,1.306);
  \fill[brick,opacity=.26] (6.540,0) rectangle (6.556,1.306);
  \fill[brick,opacity=.26] (6.556,0) rectangle (6.570,1.376);
  \fill[brick,opacity=.26] (6.570,0) rectangle (6.574,1.410);
  \fill[brick,opacity=.26] (6.574,0) rectangle (6.600,1.424);
  \fill[brick,opacity=.26] (6.600,0) rectangle (6.638,1.498);
  \fill[brick,opacity=.26] (6.638,0) rectangle (6.674,1.532);
  \fill[brick,opacity=.26] (6.674,0) rectangle (6.682,1.550);
  \fill[brick,opacity=.26] (6.682,0) rectangle (6.687,1.554);
  \fill[brick,opacity=.26] (6.687,0) rectangle (6.699,1.562);
  \fill[brick,opacity=.26] (6.699,0) rectangle (6.715,1.563);
  \fill[brick,opacity=.26] (6.715,0) rectangle (6.782,1.564);
  \fill[brick,opacity=.26] (6.782,0) rectangle (6.809,1.627);
  \fill[brick,opacity=.26] (6.809,0) rectangle (6.820,1.672);
  \fill[brick,opacity=.26] (6.820,0) rectangle (6.831,1.672);
  \fill[brick,opacity=.26] (6.831,0) rectangle (6.888,1.771);
  \fill[brick,opacity=.26] (6.888,0) rectangle (6.921,1.801);
  \fill[brick,opacity=.26] (6.921,0) rectangle (6.934,1.829);
  \fill[brick,opacity=.26] (6.934,0) rectangle (6.997,1.879);
  \fill[brick,opacity=.26] (6.997,0) rectangle (7.022,1.984);
  \fill[brick,opacity=.26] (7.022,0) rectangle (7.028,2.003);
  \fill[brick,opacity=.26] (7.028,0) rectangle (7.037,2.052);
  \fill[brick,opacity=.26] (7.037,0) rectangle (7.047,2.052);
  \fill[brick,opacity=.26] (7.047,0) rectangle (7.066,2.404);
  \fill[brick,opacity=.26] (7.066,0) rectangle (7.083,2.554);
  \fill[brick,opacity=.26] (7.083,0) rectangle (7.087,2.600);
  \fill[brick,opacity=.26] (7.087,0) rectangle (7.088,2.600);
  \fill[brick,opacity=.26] (7.088,0) rectangle (7.090,2.600);
  \fill[brick,opacity=.26] (7.090,0) rectangle (7.102,2.600);
  \fill[brick,opacity=.26] (7.102,0) rectangle (7.113,2.600);
  \fill[brick,opacity=.26] (7.113,0) rectangle (7.125,2.600);
  \fill[brick,opacity=.26] (7.125,0) rectangle (7.140,2.600);
  \fill[brick,opacity=.26] (7.140,0) rectangle (7.150,2.600);
  \draw[ink,line width=.6pt] (0.000,0.000) -- (2.860,0.000) -- (2.860,0.000) -- (2.863,0.000) -- (2.863,0.000) -- (2.867,0.000) -- (2.867,0.000) -- (2.876,0.000) -- (2.876,0.000) -- (2.888,0.000) -- (2.888,0.000) -- (2.900,0.000) -- (2.900,0.000) -- (2.924,0.000) -- (2.924,0.000) -- (2.968,0.000) -- (2.968,0.058) -- (2.999,0.058) -- (2.999,0.197) -- (3.023,0.197) -- (3.023,0.216) -- (3.059,0.216) -- (3.059,0.238) -- (3.098,0.238) -- (3.098,0.277) -- (3.103,0.277) -- (3.103,0.277) -- (3.108,0.277) -- (3.108,0.277) -- (3.113,0.277) -- (3.113,0.277) -- (3.119,0.277) -- (3.119,0.277) -- (3.125,0.277) -- (3.125,0.277) -- (3.137,0.277) -- (3.137,0.277) -- (3.177,0.277) -- (3.177,0.424) -- (3.179,0.424) -- (3.179,0.504) -- (3.180,0.504) -- (3.180,0.537) -- (3.189,0.537) -- (3.189,0.546) -- (3.212,0.546) -- (3.212,0.546) -- (3.235,0.546) -- (3.235,0.558) -- (3.344,0.558) -- (3.344,0.569) -- (3.394,0.569) -- (3.394,0.573) -- (3.412,0.573) -- (3.412,0.573) -- (6.111,0.573) -- (6.111,0.575) -- (6.136,0.575) -- (6.136,0.576) -- (6.153,0.576) -- (6.153,0.585) -- (6.184,0.585) -- (6.184,0.590) -- (6.209,0.590) -- (6.209,0.593) -- (6.226,0.593) -- (6.226,0.636) -- (6.244,0.636) -- (6.244,0.645) -- (6.266,0.645) -- (6.266,0.755) -- (6.281,0.755) -- (6.281,0.773) -- (6.288,0.773) -- (6.288,0.828) -- (6.402,0.828) -- (6.402,0.837) -- (6.410,0.837) -- (6.410,0.964) -- (6.417,0.964) -- (6.417,0.965) -- (6.441,0.965) -- (6.441,1.146) -- (6.532,1.146) -- (6.532,1.306) -- (6.540,1.306) -- (6.540,1.306) -- (6.556,1.306) -- (6.556,1.376) -- (6.570,1.376) -- (6.570,1.410) -- (6.574,1.410) -- (6.574,1.424) -- (6.600,1.424) -- (6.600,1.498) -- (6.638,1.498) -- (6.638,1.532) -- (6.674,1.532) -- (6.674,1.550) -- (6.682,1.550) -- (6.682,1.554) -- (6.687,1.554) -- (6.687,1.562) -- (6.699,1.562) -- (6.699,1.563) -- (6.715,1.563) -- (6.715,1.564) -- (6.782,1.564) -- (6.782,1.627) -- (6.809,1.627) -- (6.809,1.672) -- (6.820,1.672) -- (6.820,1.672) -- (6.831,1.672) -- (6.831,1.771) -- (6.888,1.771) -- (6.888,1.801) -- (6.921,1.801) -- (6.921,1.829) -- (6.934,1.829) -- (6.934,1.879) -- (6.997,1.879) -- (6.997,1.984) -- (7.022,1.984) -- (7.022,2.003) -- (7.028,2.003) -- (7.028,2.052) -- (7.037,2.052) -- (7.037,2.052) -- (7.047,2.052) -- (7.047,2.404) -- (7.066,2.404) -- (7.066,2.554) -- (7.083,2.554) -- (7.083,2.600) -- (7.087,2.600) -- (7.087,2.600) -- (7.088,2.600) -- (7.088,2.600) -- (7.090,2.600) -- (7.090,2.600) -- (7.102,2.600) -- (7.102,2.600) -- (7.113,2.600) -- (7.113,2.600) -- (7.125,2.600) -- (7.125,2.600) -- (7.140,2.600) -- (7.140,2.600) -- (7.150,2.600);
  \draw[rulec,line width=.6pt] (0,0) -- (7.15,0);
  \draw[rulec,line width=.6pt] (0,0) -- (0,2.85);
  \draw[rulec,line width=.6pt] (-.07,0.000) -- (0,0.000);
  \node[left,muted,font=\fontsize{5.2}{6.2}\selectfont] at (-.09,0.000) {0};
  \draw[rulec,line width=.6pt] (-.07,0.650) -- (0,0.650);
  \node[left,muted,font=\fontsize{5.2}{6.2}\selectfont] at (-.09,0.650) {350};
  \draw[rulec,line width=.6pt] (-.07,1.300) -- (0,1.300);
  \node[left,muted,font=\fontsize{5.2}{6.2}\selectfont] at (-.09,1.300) {700};
  \draw[rulec,line width=.6pt] (-.07,1.950) -- (0,1.950);
  \node[left,muted,font=\fontsize{5.2}{6.2}\selectfont] at (-.09,1.950) {1050};
  \draw[rulec,line width=.6pt] (-.07,2.600) -- (0,2.600);
  \node[left,muted,font=\fontsize{5.2}{6.2}\selectfont] at (-.09,2.600) {1400};
  \node[below,muted,font=\fontsize{5.2}{6.2}\selectfont] at (0.000,-.04) {0};
  \node[below,muted,font=\fontsize{5.2}{6.2}\selectfont] at (1.715,-.04) {25};
  \node[below,muted,font=\fontsize{5.2}{6.2}\selectfont] at (3.430,-.04) {50};
  \node[below,muted,font=\fontsize{5.2}{6.2}\selectfont] at (5.144,-.04) {75};
  \node[below,muted,font=\fontsize{5.2}{6.2}\selectfont] at (6.859,-.04) {100};
  \draw[ink,line width=.35pt,dashed] (3.440,0) -- (3.440,2.85);
  \node[right,ink,font=\fontsize{5.6}{6.7}\selectfont\bfseries] at (3.520,2.75) {t1};
  \node[right,brick,font=\fontsize{5.2}{6.2}\selectfont] at (3.520,2.52) {50 GW};
  \draw[ink,line width=.35pt,dashed] (6.387,0) -- (6.387,2.85);
  \node[left,ink,font=\fontsize{5.6}{6.7}\selectfont\bfseries] at (6.307,2.75) {t2};
  \node[left,brick,font=\fontsize{5.2}{6.2}\selectfont] at (6.307,2.52) {93 GW};
  \draw[brick,line width=.7pt] (0,0.573) -- (7.15,0.573);
  \node[right,fill=brick,text=white,inner sep=1pt,font=\fontsize{5.0}{6.0}\selectfont] at (.06,0.703) {t1 $\rightarrow$ 309};
  \draw[brick,line width=.7pt] (0,0.828) -- (7.15,0.828);
  \node[right,fill=brick,text=white,inner sep=1pt,font=\fontsize{5.0}{6.0}\selectfont] at (.06,0.958) {t2 $\rightarrow$ 446};
\end{tikzpicture}}

    \vspace{0.5em}
    \raggedright
    {\fontsize{5.9}{7.4}\selectfont\color{muted}
     \textcolor{muted!60}{\rule[0ex]{0.6em}{0.6em}}~inflexible hydro\hfill
     \textcolor{tealx!45}{\rule[0ex]{0.6em}{0.6em}}~flexible hydro\hfill
     \textcolor{brick!40}{\rule[0ex]{0.6em}{0.6em}}~thermal by cost\par
     \textcolor{ink}{\rule[0.22ex]{0.9em}{0.4pt}}~quantity of the hour\hfill
     \textcolor{brick}{\rule[0.22ex]{0.9em}{0.8pt}}~model price\hfill\mbox{}\par}

    \vspace{0.4em}
    \gloss{The dashed verticals are quantities, what the fleet had to deliver in each
    hour; the horizontals are prices. They pull apart when flexible hydro can no
    longer reach the peak. Ramp of 14.3\%/h of flexible hydro, the day's own.}
  \end{column}

  \begin{column}{0.425\textwidth}
    \vspace{0.5em}
    {\fontsize{6.5}{8.4}\selectfont
    \setlength{\tabcolsep}{1.9pt}
    \renewcommand{\arraystretch}{1.12}
    \begin{tabular}{l r r r @{\hspace{0.55em}} r r r}
      \toprule
      {\fontsize{5.6}{6}\selectfont R\$ m}
        & \multicolumn{3}{c}{\textbf{Without bat. \ 309\,$\rightarrow$\,446}}
        & \multicolumn{3}{c}{\textbf{With bat. \ 309\,$\rightarrow$\,347}} \\
      \cmidrule(r){2-4}\cmidrule(l){5-7}
       & Rev. & Cost & Profit & Rev. & Cost & Profit \\
      \midrule
      Inflex. hydro & 31.5 & 0.0 & 31.5 & 27.3 & 0.0 & 27.3 \\
      Flexible hydro & 17.7 & 12.3 & 5.4 & 14.4 & 12.9 & 1.5 \\
      Thermal        & 7.9 & 4.3 & 3.6 & 6.0 & 3.4 & 2.6 \\
      Battery        & & & & \hg{0.1} & \hg{0.0} & \hg{0.1} \\
      \addlinespace[0.1em]\midrule
      Net load       & \hl{$-$57.0} & & & \hl{$-$47.8} & & \\
      \addlinespace[0.1em]\midrule
      \textbf{Total} & \textbf{0.0} & \textbf{16.5} & \textbf{40.5}
                       & \textbf{0.0} & \textbf{16.3} & \textbf{31.5} \\
      \bottomrule
    \end{tabular}}

    \vspace{0.35em}
    \gloss{With 2,000 MWh of battery. Revenue sums to zero in both cases.}

    \vspace{0.4em}
    \thesis{\fontsize{7.1}{9}\selectfont
      Demand pays R\$ 9.2 million less, and operating cost falls only R\$ 0.2
      million: most of it is a transfer from inframarginal generators to
      consumers, and the battery keeps R\$ 0.1 million.}
  \end{column}
  \end{columns}
\end{frame}
